% Generated draft: requires mathematical and source review.
\documentclass[11pt]{article}
\usepackage[margin=25mm]{geometry}
\usepackage{fontspec}
\setmainfont{Latin Modern Roman}
\usepackage{amsmath,amssymb,mathtools,unicode-math}
\setmathfont{Latin Modern Math}
\usepackage{enumitem,needspace,xurl,hyperref}
\setlength{\parindent}{0pt}
\setlength{\parskip}{5pt}
\setlength{\emergencystretch}{3em}
\newcommand{\N}{\mathbb{N}}
\newcommand{\Z}{\mathbb{Z}}
\newcommand{\Q}{\mathbb{Q}}
\newcommand{\R}{\mathbb{R}}
\newcommand{\C}{\mathbb{C}}
\newcommand{\F}{\mathbb{F}}
\newcommand{\A}{\mathbb{A}}
\newcommand{\PP}{\mathbb{P}}
\newcommand{\E}{\mathbb{E}}
\newcommand{\eps}{\varepsilon}
\newcommand{\dd}{\,\mathrm{d}}
\begin{document}
\section*{UnsolvedMath ID: 30001473}
\textbf{Source problem code:} catalogue entry
\subsection{Definitions and mathematical statement}
A Coxeter group is a group admitting a presentation
\[
W=\left\langle S\ \middle|\ (st)^{m_{st}}=1\text{ whenever }m_{st}<\infty\right\rangle,
\]
where \(m_{ss}=1\) and \(m_{st}=m_{ts}\in\{2,3,\ldots\}\cup\{\infty\}\) for distinct \(s,t\in S\); an exponent \(\infty\) imposes no relation. The supplied statement specifies no restriction on the cardinality of \(S\).

One standard group-theoretic meaning of \emph{large} is the existence of a subgroup \(W_0\leq W\) of finite index and a surjective homomorphism \(W_0\twoheadrightarrow F_2\), where \(F_2\) is the free group on two generators. The supplied record does not define \emph{large}, so use of this meaning is an explicit, unverified interpretation.

A group \(H\) is word-hyperbolic if it has a finite generating set for which the Cayley graph, with its graph metric, has \(\delta\)-thin geodesic triangles for some \(\delta\geq0\): every side of each geodesic triangle lies in the closed \(\delta\)-neighborhood of the union of the other two sides. A homomorphism \(\varphi:W\to H\) is surjective if \(\varphi(W)=H\).

The supplied question is: \emph{Does every large Coxeter group surject onto a hyperbolic group?} Under the interpretations above, its literal quantified reading is
\[
\begin{gathered}
\forall W\text{ a Coxeter group},\\
\left(\exists W_0\leq W:\ [W:W_0]<\infty\text{ and }W_0\twoheadrightarrow F_2\right)
\ \Longrightarrow\ 
\left(\exists H\text{ word-hyperbolic},\ \exists\varphi:W\twoheadrightarrow H\right).
\end{gathered}
\]
\textbf{Unresolved qualifications.} Word-hyperbolicity in its standard sense permits finite groups and virtually cyclic groups. The record does not state whether these targets are permitted, or whether the intended target must instead be infinite or non-elementary. Here \emph{non-elementary} would mean neither finite nor virtually cyclic, and \emph{virtually cyclic} means having a cyclic subgroup of finite index. No such additional target condition is inserted into the displayed reading. The intended meaning of \emph{large}, any finite-rank hypothesis, and the intended restriction on the quotient require source clarification before this can be treated as a precise substantive open-problem formulation.
\subsection{Short English statement}
The supplied question asks whether every large Coxeter group has a hyperbolic quotient. A literal reading is displayed using an explicitly qualified standard interpretation of largeness. The source record leaves the meaning of largeness, any rank restriction, and possible exclusions of finite or virtually cyclic quotients unresolved. Its supplied open status is dated curation, not an independently verified conclusion.
\subsection{Sources}
\begin{itemize}
\item Local registry: \texttt{data/unsolvedmath/problems.json}, permanent UnsolvedMath ID \(30001473\), source problem code \texttt{catalogue entry}, title \emph{Hyperbolic Quotients of Large Coxeter Groups}; collection: \emph{Oberwolfach Reports Open Problems}.
\item Supplied source citation: \emph{Non-positive Curvature and Geometric Structures in Group Theory} (2010), Oberwolfach Reports, pp.~1165--1224, DOI \texttt{10.4171/owr/2010/20}. Supplied source URL and literature-source link labeled ``OWR 2010/20'': \url{https://doi.org/10.4171/owr/2010/20}. The record describes these reports as documenting mathematical workshops and their open-problem discussions. No individual proposer is supplied.
\item Supplied literature assessment, checked \textbf{2026-08-21}: status \emph{open}. The supplied assessment says that the source explicitly poses the surjectivity question and that its literature search found no general resolution. This is supplied dated curation derived from the OWR record, not independently verified source consultation or a new literature assessment.
\item The registry labels the statement ``exact'' and reports that manual review found the extracted statement self-contained, mathematically coherent, and free of LaTeX or report-layout artifacts. These are supplied verification claims; they do not resolve the formulation qualifications identified here.
\end{itemize}
\end{document}
